\scanpage{192}

$\underline{\overline{\mathrm{III}}}$\quad $G=\underset{\substack{\parallel\\F}}{k}\mathbin{.}\underset{\substack{\parallel\\H}}{k^*}$ produit semi-direct, $k^*$ opérant dans $k$ par $\sigma_p(\lambda)t=\lambda^pt$, ($n$ entier rationnel \uline{non nul}).

On suppose $k$ de car. 0. L'algèbre de Lie \deleted{\unsure{de}} $\unsure{\mathfrak g|\mathscr G}$ de $G$ est extension de $\unsure{\mathscr F|\mathfrak F|F}\simeq k$ par $\unsure{\mathscr H|\mathfrak H|H}\simeq k$, $\mathscr H$ opérant sur $\mathscr F$ \deleted{par} ainsi ($X$ et $Y$ sont les \unsure{éléments-base|éléments-bases} canoniques de $\mathscr F$ et $\mathscr H$) :

\[[Y,X]=pX\tag*{(1)}\]

Considérons une représentation linéaire indécomposable \unsure{rationnelle|rationelle} $\varphi$ de $\deleted{\unsure{G_2(\mathscr F)}}G$, soit $\varphi'$ sa différentielle. On a (décomposition en sous-espaces propres \unsure{relt|rel\textsuperscript{ts}} à $H\subset G$)

\[V=\sum V_n\]

($Y.a=\deleted{p}na$ si $a\in V_n$). La condition (1) montre que les composantes $X_{\deleted{\unsure{ij|i}}ji}\in L(V_i,V_j)$ de $X$ satisfont :

\[(j-i)X_{ji}=pX_{ji}\]

d'où $X_{ji}=0$ sauf si $j-i=p$.

% archive sheet number: 191
